% =============================================================================
% 04_closing.tex — Slides S01-067 to S01-070
% Session Closing, Master Synthesis, 5 Engineering Rules & Bridge to Session 02
% =============================================================================

% -----------------------------------------------------------------------------
% Slide ID: S01-067
% Narrative Role: Master Synthesis
% Cognitive Objective: Unify the 3 session blocks into one connected mental model.
% Plan Source: slide_plan.md / S01-067
% -----------------------------------------------------------------------------
\begin{frame}{Session 01 Master Synthesis: The Unified Mental Model}
  \begin{columns}[T,onlytextwidth]
    \begin{column}{0.31\textwidth}
      \begin{tcolorbox}[enhanced,arc=1.4mm,boxrule=.4pt,colback=UpGradMist,colframe=UpGradRed,
        borderline north={2.5pt}{0pt}{UpGradRed},title={\bfseries\scriptsize Tier 1: Dynamic Flexibility},coltitle=UpGradNight]
        \scriptsize
        \textbf{Stack Pointers \& Heap}\\[2pt]
        Dynamic typing, closures, decorators, and expressive metaprogramming via boxed \InlineCode{PyObject} structs.
      \end{tcolorbox}
    \end{column}
    \begin{column}{0.31\textwidth}
      \begin{tcolorbox}[enhanced,arc=1.4mm,boxrule=.4pt,colback=UpGradMist,colframe=AWSBlue,
        borderline north={2.5pt}{0pt}{AWSBlue},title={\bfseries\scriptsize Tier 2: Stream Efficiency},coltitle=UpGradNight]
        \scriptsize
        \textbf{Hash Tables \& Generators}\\[2pt]
        $O(1)$ fast lookups with sets/dicts; $O(1)$ constant memory streaming with lazy \InlineCode{yield} iterators.
      \end{tcolorbox}
    \end{column}
    \begin{column}{0.31\textwidth}
      \begin{tcolorbox}[enhanced,arc=1.4mm,boxrule=.4pt,colback=UpGradMist,colframe=AWSEmerald,
        borderline north={2.5pt}{0pt}{AWSEmerald},title={\bfseries\scriptsize Tier 3: Hardware Speed},coltitle=UpGradNight]
        \scriptsize
        \textbf{Contiguous NumPy Arrays}\\[2pt]
        Zero-copy strided views, SIMD vector registers, and zero-memory broadcasting engine.
      \end{tcolorbox}
    \end{column}
  \end{columns}
  \vspace{0.18cm}
  \TakeawayBanner{Python provides the complete spectrum: high-level expressiveness down to native hardware C-speed.}
\end{frame}

% -----------------------------------------------------------------------------
% Slide ID: S01-068
% Narrative Role: North-Star Resolution
% Cognitive Objective: Deliver a definitive, rigorous answer to the opening question.
% Plan Source: slide_plan.md / S01-068
% -----------------------------------------------------------------------------
\begin{frame}{Resolution of the North-Star Question}
  \begin{center}
    \vspace{0.20cm}
    \begin{tcolorbox}[enhanced,arc=2mm,boxrule=.6pt,colback=UpGradMist,colframe=UpGradRule,
      borderline west={3.5pt}{0pt}{AWSEmerald},left=5mm,right=5mm,top=4mm,bottom=4mm,width=0.94\linewidth]
      {\small\bfseries\color{UpGradNight}
        “Python achieves its dynamic flexibility because variables are named pointers to boxed heap objects (\InlineCode{PyObject}). This indirection creates severe overhead in tight loops through pointer chasing, boxed headers, and dynamic dispatch.\\[6pt]
        NumPy breaks through this speed wall by bypassing Python objects entirely, executing compiled SIMD operations directly on contiguous, strided memory buffers.”
      }
    \end{tcolorbox}
  \end{center}
  \vfill
  \TakeawayBanner{We have resolved the paradox: Flexibility in Python's high layer; raw C performance in NumPy's buffer layer.}
\end{frame}

% -----------------------------------------------------------------------------
% Slide ID: S01-069
% Narrative Role: Engineering Takeaways / Session Summary
% Cognitive Objective: Provide an actionable summary checklist for production Python & ML engineering.
% Plan Source: slide_plan.md / S01-069
% -----------------------------------------------------------------------------
\begin{frame}{Session 01 Summary: 5 Golden Engineering Rules}
  \begin{tcolorbox}[enhanced,arc=1.4mm,boxrule=.4pt,colback=UpGradMist,colframe=UpGradRule,left=3.5mm,right=3.5mm,top=1.8mm,bottom=1.8mm]
    \footnotesize
    \begin{enumerate}\setlength{\itemsep}{0.05cm}
      \item \textbf{Never use mutable default arguments:} Always use immutable sentinels (\InlineCode{def f(x=None)}).
      \item \textbf{Beware of shallow copies on nested data:} Use \InlineCode{copy.deepcopy()} when modifying compound graphs.
      \item \textbf{Preserve decorator metadata:} Always wrap decorator wrappers with \InlineCode{@functools.wraps(func)}.
      \item \textbf{Stream large datasets with generators:} Use \InlineCode{yield} to maintain an $O(1)$ constant memory footprint.
      \item \textbf{Never write \InlineCode{for} loops for array math:} Express calculations as vectorized NumPy ufuncs and broadcasts.
    \end{enumerate}
  \end{tcolorbox}
  \vfill
  \TakeawayBanner{Adhering to these 5 rules eliminates 99\% of production memory leaks and performance bottlenecks.}
\end{frame}

% -----------------------------------------------------------------------------
% Slide ID: S01-070
% Narrative Role: Course Bridge / Outro
% Cognitive Objective: Spark curiosity for Session 02 (Applied Linear Algebra via NumPy).
% Plan Source: slide_plan.md / S01-070
% -----------------------------------------------------------------------------
\begin{frame}{Bridge to Session 02: Applied Linear Algebra via NumPy}
  \begin{center}
    \vspace{0.05cm}
    \begin{tcolorbox}[enhanced,arc=2mm,boxrule=.6pt,colback=UpGradNight,colframe=AWSSquid,
      left=5mm,right=5mm,top=2.5mm,bottom=2.5mm,width=0.92\linewidth]
      {\UpGradDisplayFont\fontsize{10.5}{14}\selectfont\bfseries\color{white}
        “We now know how to compute on contiguous multi-dimensional arrays at native hardware speed. But what do these arrays, matrices, and transformations actually mean geometrically when building machine learning algorithms?”
      }
    \end{tcolorbox}
  \end{center}
  \vspace{0.10cm}
  \begin{columns}[c,onlytextwidth]
    \begin{column}{0.31\textwidth}
      \StatCard{AWSOrange}{Session 02}{Vector Spaces}{Data points as geometric arrows}
    \end{column}
    \begin{column}{0.31\textwidth}
      \StatCard{AWSBlue}{Matrix Math}{Linear Maps}{Coordinate space transformations}
    \end{column}
    \begin{column}{0.31\textwidth}
      \StatCard{AWSEmerald}{SVD \& PCA}{Eigendecomp}{The master data compression tool}
    \end{column}
  \end{columns}
\end{frame}
